%% gte_search_targets_table.tex — includable fragment for P00, P48, etc.
%% Source: GTE novel-particle register (P02) and dark-sector complement (P29).

%% ---- colours --------------------------------------------------------
\definecolor{tier1}{rgb}{0.0,0.40,0.0}
\definecolor{tier2}{rgb}{0.55,0.27,0.07}
\definecolor{tier3}{rgb}{0.50,0.00,0.50}
\definecolor{darkgray}{gray}{0.35}

%% ---- column types (local to this fragment) ------------------------
\newcolumntype{L}[1]{>{\raggedright\arraybackslash}p{#1}}
\newcolumntype{C}[1]{>{\centering\arraybackslash}p{#1}}

%% ---- macros ---------------------------------------------------------
\newcommand{\tOne}{\textcolor{tier1}{\textbf{Tier~1}}}
\newcommand{\tTwo}{\textcolor{tier2}{\textbf{Tier~2}}}
\newcommand{\tThree}{\textcolor{tier3}{\textbf{Tier~3}}}
\newcommand{\MeV}{\,\mathrm{MeV}}
\newcommand{\GeV}{\,\mathrm{GeV}}

\subsection*{Overview}

This table summarizes novel particle-mass predictions of the Universal Generative
Principle (UGP) framework that may be testable against existing or near-term
accelerator data.
Items GTE-P7 through GTE-P11 are a subset of the full novel-particle register in
P02~\cite{SpivackGTESpectrum}; P29 provides the dark-sector complement~\cite{SpivackMirrorBraidAtlas}.
A broader derivation context (Standard Model parameters, gauge couplings, fermion
masses) is in P01~\cite{SpivackSM}.

\subsection*{Key Signature Property}

All 9 genuinely novel candidates from the GTE spectrum survey are predicted
\textbf{stable} (lifetime $\tau \geq 10^{30}$~s).  This means the
search strategy differs from ordinary resonance hunting:

\begin{itemize}[leftmargin=2em,itemsep=2pt]
  \item Stable neutral states ($Q=0$) appear as \textbf{missing energy/momentum}.
  \item Stable charged states appear as \textbf{anomalous long-lived tracks} or
        \textbf{heavy stable charged particles} (HSCP).
  \item Standard resonance-width fits in hadronic spectra would \emph{not} catch
        these; one instead looks for a \emph{sharp excess} at a specific invariant
        mass, or an anomalous \emph{endpoint feature}, in an otherwise smooth
        background.
\end{itemize}

The highest-priority state (\textbf{GTE-P7}) is electrically neutral ($Q=0$,
color-singlet, Lean~4 certified with zero axioms).  It is \emph{not} a dark
photon: the framework predicts $\varepsilon=0$ exactly (no kinetic mixing), so
existing dark-photon exclusions do \emph{not} apply to GTE-P7.

\bigskip

\subsection*{Predicted Targets by Search Tier}

\noindent
\textcolor{tier1}{\rule{1em}{0.8em}}~Tier~1 = testable with existing archival data now.\quad
\textcolor{tier2}{\rule{1em}{0.8em}}~Tier~2 = Belle~II / LHCb Run~3 ($\sim$2026--2028).\quad
\textcolor{tier3}{\rule{1em}{0.8em}}~Tier~3 = requires dedicated low-energy experiment.

\medskip

\setlength{\tabcolsep}{5pt}
\renewcommand{\arraystretch}{1.50}
\small
\begin{longtable}{%
  C{0.060\textwidth}%
  C{0.095\textwidth}%
  C{0.065\textwidth}%
  C{0.115\textwidth}%
  L{0.240\textwidth}%
  L{0.270\textwidth}%
}
\toprule
\textbf{State} & \textbf{Mass} & \textbf{Unc.} & \textbf{QN} &
\textbf{Search channel / experiment} & \textbf{Notes} \\
\midrule
\endfirsthead
\multicolumn{6}{l}{\footnotesize\itshape (continued)} \\
\toprule
\textbf{State} & \textbf{Mass} & \textbf{Unc.} & \textbf{QN} &
\textbf{Search channel / experiment} & \textbf{Notes} \\
\midrule
\endhead
\bottomrule
\endlastfoot

%% ----- Tier 1 -----
\multicolumn{6}{l}{\tOne} \\[2pt]
\midrule

\textbf{GTE-P7}
  & $211.9\MeV$
  & $\pm14\%$
  & $Q=0$,\newline singlet,\newline $s=\tfrac{1}{2}$\newline Lean-cert.
  & BaBar archival: dimuon threshold / $e^+e^-\to\gamma X$\newline
    LHCb archival: dimuon at $2m_\mu$\newline
    Belle~II: mono-photon $M_\mathrm{recoil}\approx212\MeV$
  & Sits at dimuon threshold ($2m_\mu=211.4\MeV$).
    Signal = missing invariant mass recoiling against photon.
    \emph{Not} a dark photon; existing $\varepsilon$-exclusions do not apply.
    50~ab$^{-1}$ at Belle~II constrains $\varepsilon<10^{-3}$.
    \textbf{Highest near-term priority.} \\[4pt]

\textbf{GTE-P8}
  & $298\MeV$
  & $\pm14\%$
  & heuristic: quark-like\newline(Möbius product $= -1$)
  & BESIII existing data: $\pi\pi\eta$ and $2\gamma$ spectra\newline
    BaBar/Belle archival: same channels
  & Look for a \emph{stable} peak near 298~MeV not matching any
    known resonance.  No known PDG state within 10\% of this mass. \\[4pt]

\textbf{GTE-P9}
  & $561\MeV$
  & $\pm20\%$
  & heuristic:\newline sector indet.
  & BESIII: $J/\psi$ radiative decays, 550--575~MeV window\newline
    LHCb archival: inclusive spectra
  & Between $\eta'(958)$ and $\phi(1020)$; look for a stable state
    distinct from both. Window: $\sim449$--$673\MeV$. \\[8pt]

%% ----- Tier 2 -----
\multicolumn{6}{l}{\tTwo} \\[2pt]
\midrule

\textbf{GTE-P3}
  & $\sim800\MeV$\newline (band: 796--850)
  & $\pm25\%$
  & charge uncertain
  & Belle~II, LHCb Run~3:\newline $\pi\pi K\bar{K}$ mass spectra
  & In the $f_0(980)/\omega(782)$ region.  Look for a \emph{stable}
    charged or neutral state not matching $f_0$, $\omega$, or $\rho$.
    Window: $\sim600$--$1063\MeV$. \\[4pt]

\textbf{GTE-P2}
  & $107\MeV$\newline (cluster: 107--110)
  & $\pm14\%$
  & charge uncertain
  & NA62 rare decays\newline
    DarkLight (JLab)\newline
    Long-lived / displaced-vertex searches
  & 1.75~MeV above the muon; a \emph{distinct} state (different
    GTE triple: $n=4702$ vs muon $n=42$).  Window: $\sim92$--$125\MeV$. \\[4pt]

\textbf{GTE-P6}
  & $30.9\MeV$\newline {\small(29--33)}
  & $\pm14\%$
  & charge uncertain
  & NA62\newline FASER at LHC\newline Dedicated dark-sector fixed-target
  & Window: $\sim26$--$37\MeV$.  No known PDG resonance in range. \\[4pt]

\textbf{GTE-P10}
  & $\sim1.1$--$1.35\GeV$\newline (735 cluster members)
  & $\pm20\%$
  & charge uncertain
  & LHCb Run~3\newline BESIII charm threshold
  & Charm-adjacent; internally consistent with P01 first-principles
    charm prediction (1276~MeV) from an independent derivation path.
    Window: $\sim880$--$1620\MeV$. \\[4pt]

\textbf{GTE-P11}
  & $\sim1.6$--$1.9\GeV$\newline (1919 cluster members)
  & $\pm25\%$
  & charge uncertain
  & LHCb Run~3\newline BESIII
  & Tau-adjacent; internally consistent with P01 tau prediction
    (1775~MeV).  Window: $\sim1.2$--$2.4\GeV$. \\[8pt]

%% ----- Tier 3 -----
\multicolumn{6}{l}{\tThree} \\[2pt]
\midrule

\textbf{GTE-P1}
  & $2.97\MeV$
  & $\pm1.4\%$
  & lepton-like\newline ($\mu$-par.\ $+1$)
  & DAFNE (Frascati)\newline
    Dedicated MeV-scale experiments
  & Sharpest mass prediction in the entire million-candidate dataset.
    Below most collider thresholds.  Window: $2.93$--$3.01\MeV$. \\[8pt]

\bottomrule
\end{longtable}
\normalsize
\setlength{\tabcolsep}{6pt}

\bigskip

\subsection*{Dark Sector Mirror Branch (from P29)}

These are three \emph{dark singlet leptons} ($Q_\mathrm{EM}=0$, no weak charge,
color-singlet) derived from the GTE mirror branch with zero free parameters.
Mass ratios G2/G1 $\approx 45.3$ and G3/G2 $\approx 147.3$ are fixed structural
predictions, falsifiable if states are found at different mass ratios.

\medskip
\setlength{\tabcolsep}{5pt}
\renewcommand{\arraystretch}{1.4}
\small
\begin{tabular}{L{0.24\textwidth}C{0.10\textwidth}L{0.13\textwidth}L{0.40\textwidth}}
\toprule
\textbf{State} & \textbf{Mass} & \textbf{Certification} & \textbf{Search / comment} \\
\midrule
$\chi_1$ (dark lepton G1) & $0.5406\MeV$ & GTE arithmetic & Sub-MeV DM via Higgs portal;
  below CRESST-III sensitivity; CMB spectral distortions (future) \\
$\chi_2$ (dark lepton G2) & $24.47\MeV$  & GTE arithmetic & MeV-scale Higgs-portal DM;
  two-peak signature with $\chi_1$\allowbreak\ at fixed ratio 45.3 \\
$\chi_3$ (dark lepton G3) & $3604.68\MeV$ & GTE arithmetic & LHC Higgs invisible width;
  mono-jet; coupling suppressed by\newline $\lambda_s \sim 10^{-6}$ \\
Dark photon $A'$ & \textbf{does not exist} & Lean-cert.\ ($\varepsilon=0$) &
  Null results from NA64/BaBar/Belle~II are \emph{consistent} with the framework \\
\bottomrule
\end{tabular}
\normalsize
\setlength{\tabcolsep}{6pt}

\bigskip

\subsection*{Priority Guidance for Existing Data}

If access is limited to archival data from one experiment, the recommended search
order is:

\begin{enumerate}[leftmargin=2.5em,itemsep=3pt]
  \item \textbf{GTE-P7 at $\approx212\MeV$} in BaBar or LHCb archival dimuon/missing-energy
    data.  The signature is a \emph{stable, neutral} state appearing as missing invariant
    mass recoiling against an ISR photon.  BaBar has published mono-photon searches in
    this mass range; a re-analysis asking specifically for a \emph{stable} (not radiatively
    decaying) state at the dimuon threshold would be the most targeted test.

  \item \textbf{GTE-P8 at $\approx298\MeV$} in BESIII hadronic data.  Look for an anomalous
    bump in the $\pi\pi\eta$ or $\gamma\gamma$ invariant mass spectrum near 298~MeV,
    distinct from the known $\eta(548)$, $\eta'(958)$, and $\omega(782)$ peaks.

  \item \textbf{GTE-P9 at $\approx561\MeV$} in BESIII $J/\psi$ radiative decays.  The
    550--575~MeV window is relatively clean of established resonances and has been
    surveyed for other purposes; a targeted stable-state search is straightforward.

  \item \textbf{GTE-P10/P11 (1.1--1.9~GeV band)} in LHCb or BESIII prompt data:
    anomalously long-lived or stable new states in the charm/tau mass region with
    unexpected quantum numbers.
\end{enumerate}

\medskip

\subsubsection*{A note on search strategy}

Because all GTE novel candidates are predicted stable, the standard procedure of
looking for a Breit--Wigner peak with a measured width does not apply.  The
relevant search is instead for:
\begin{itemize}[leftmargin=2em,itemsep=2pt]
  \item A \emph{narrow excess} (width consistent with detector resolution) at a
    specific invariant mass, \emph{with no matching known SM resonance}.
  \item For neutral states: an \emph{endpoint or threshold feature} in
    missing-energy spectra.
  \item For charged states: an HSCP (heavy stable charged particle) signature ---
    anomalously high $dE/dx$ or anomalous time-of-flight.
\end{itemize}

Absence of any signal in the mass window listed above \emph{falsifies} the
corresponding GTE prediction.  The framework makes crisp falsification predictions:
each Tier~1 target can be definitively excluded or confirmed with data already
on tape.

\bigskip

\subsection*{Oscillation-Parameter Falsification Target (JUNO / NuFIT)}

The GTE seesaw mechanism predicts a specific neutrino mass ordering testable by
reactor-oscillation experiments.  Unlike the particle-mass targets above, this
prediction is an oscillation-parameter observable, not a new-state search.

\medskip
\setlength{\tabcolsep}{5pt}
\renewcommand{\arraystretch}{1.4}
\small
\begin{tabular}{L{0.19\textwidth}L{0.27\textwidth}C{0.07\textwidth}L{0.14\textwidth}L{0.27\textwidth}}
\toprule
\textbf{Observable} & \textbf{GTE Prediction} & \textbf{Cert.} &
  \textbf{Experiment} & \textbf{Falsifying condition} \\
\midrule
Neutrino mass ordering &
  Normal Ordering (NH) forced by seed ordering $b_{R,1}{<}b_{R,2}{<}b_{R,3}$
  via seesaw $m_{\nu,k}=C\!\cdot\!b_{R,k}^{29/9}$; masses
  $\{m_1,m_2,m_3\}\approx\{0.68,\;8.6,\;50.1\}$~meV;
  $\Delta m^2_{21}=7.37{\times}10^{-5}$~eV$^2$,
  $\Delta m^2_{31}=2.51{\times}10^{-3}$~eV$^2$ (both within $1\%$ of NuFIT~6.0 IC24 NH) &
  \CatAL &
  JUNO reactor $\bar{\nu}_e$ disappearance; NuFIT global fits &
  Confirmed Inverted Hierarchy (IH) at ${>}3\sigma$: would falsify the
  seed-ordering mechanism of P21 and the octonionic triality confirmation of P55 \\
\bottomrule
\end{tabular}
\normalsize
\setlength{\tabcolsep}{6pt}

\medskip

\noindent\textit{Source papers:} P21~\cite{SpivackNeutrino} (primary seesaw derivation);
P55~\cite{SpivackOctonionicShadow} (octonionic triality corroboration,
\lean{seesaw\_normal\_ordering\_from\_seed\_ordering}, zero sorry).
Current NuFIT~6.0 IC24 NH data disfavour IH at $\Delta\chi^2=6.1$ (${\approx}2.5\sigma$),
consistent with the GTE prediction.
